\documentclass[11pt]{article}

\usepackage[margin=1in]{geometry}
\usepackage[T1]{fontenc}
\usepackage[utf8]{inputenc}
\usepackage{lmodern}
\usepackage{amsmath,amssymb,mathtools,bm}
\usepackage{booktabs}
\usepackage{array}
\usepackage{xcolor}
\usepackage{enumitem}
\usepackage[hidelinks]{hyperref}
\usepackage{listings}
\usepackage{longtable}
\usepackage{caption}

\definecolor{codegray}{RGB}{245,245,245}
\definecolor{codeborder}{RGB}{210,210,210}
\definecolor{darkblue}{RGB}{20,75,140}

\lstdefinelanguage{Julia}{
  morekeywords={function,struct,mutable,abstract,type,where,return,for,in,end,if,else,elseif,while,begin,let,do,try,catch,finally,module,using,import,export,const,global,local,break,continue,true,false,Nothing,Union,Vector,Matrix,Array,ComplexF64,Float64,Int,Bool,Symbol,Dict,Tuple,NamedTuple},
  sensitive=true,
  morecomment=[l]{\#},
  morestring=[b]",
}
\lstset{
  language=Julia,
  backgroundcolor=\color{codegray},
  frame=single,
  rulecolor=\color{codeborder},
  basicstyle=\ttfamily\small,
  keywordstyle=\color{darkblue}\bfseries,
  commentstyle=\color{gray},
  stringstyle=\color{teal!70!black},
  columns=fullflexible,
  keepspaces=true,
  showstringspaces=false,
  breaklines=true,
  tabsize=2
}

\newcommand{\C}{\mathbb{C}}
\newcommand{\R}{\mathbb{R}}
\newcommand{\F}{\mathbb{F}}
\newcommand{\E}{\mathbb{E}}
\newcommand{\dd}{\mathrm{d}}
\newcommand{\jj}{\mathrm{j}}
\newcommand{\norm}[1]{\left\lVert #1 \right\rVert}
\newcommand{\abs}[1]{\left\lvert #1 \right\rvert}
\newcommand{\re}{\operatorname{Re}}
\newcommand{\im}{\operatorname{Im}}
\newcommand{\sech}{\operatorname{sech}}
\newcommand{\diag}{\operatorname{diag}}
\newcommand{\clip}{\operatorname{clip}}

\title{SC-JUNA Manual Gradient Implementation Plan\\
\large Hand-Derived Real Gradients for LDPC-Coupled Receiver Updates}
\author{Implementation specification for Julia}
\date{\today}

\begin{document}
\maketitle

\begin{abstract}
This document specifies how to implement the SC-JUNA bit-latent update without relying on automatic differentiation.  The core decision is simple: use closed-form least-squares updates for the branch model and combiner, and use a hand-derived real gradient for the latent vector $\bm z\in\R^{n_b}$.  Automatic differentiation may be used only as a test-time gradient checker.  Production code should use the formulas in this document.
\end{abstract}

\tableofcontents
\clearpage

\section{Implementation decision}

The implementation should \emph{not} require automatic differentiation for the main receiver loop.  SC-JUNA has a block structure:
\begin{align}
\mathbf C &\leftarrow \text{ridge least squares},\\
\mathbf W &\leftarrow \text{ridge least squares or block-tridiagonal least squares},\\
\bm z &\leftarrow \text{manual real-gradient update}.
\end{align}
The only nontrivial gradient is the gradient with respect to the real bit-latent vector $\bm z$.  That gradient is short, stable, and easy to unit-test by finite differences.

\paragraph{Production rule.}
Use the hand-derived gradient in Section~\ref{sec:z-gradient}.  Do not use complex automatic differentiation inside the production solver.

\paragraph{Testing rule.}
Use central finite differences as the required gradient check.  AD may be used as an optional secondary check, but the finite-difference test is the authoritative test.

\paragraph{Reason.}
The $\bm z$ variables are real.  The observation residuals are complex.  Several Julia AD packages handle complex-valued intermediate computations differently.  A manual real gradient avoids package-specific complex-AD conventions and makes the implementation deterministic.

\section{Version-1 restrictions}

To keep the implementation deterministic, version 1 uses the following restrictions.

\begin{table}[h]
\centering
\caption{Version-1 choices.  These are implementation constraints, not theoretical limits.}
\begin{tabular}{@{}llp{0.45\linewidth}@{}}
\toprule
Item & Value & Reason \\
\midrule
Noise weight & $\Sigma_k^{-1}=\sigma^{-2}I_M$ & Gives simple closed-form $\mathbf C$ update and simple gradients. \\
Partial-FFT windows & nested rectangular only & Guarantees partition-of-unity and supports branch merging. \\
Latents & real $z_i\in\R$ only & Avoids complex AD and makes parity penalty simple. \\
Symbol map & relaxed QPSK, $s=(\xi_I+\jj\xi_Q)/\sqrt2$ & Smooth map from bits to complex symbols. \\
Neighborhood edge handling & skip invalid neighbors & No wraparound at band edges. \\
Optimizer for $\bm z$ & Adam or normalized gradient descent & Both use the same manual gradient. \\
AD & test-only & Never required for production. \\
\bottomrule
\end{tabular}
\end{table}

\section{Notation and array layout}

All Julia arrays use 1-based indexing.

\begin{longtable}{@{}p{0.23\linewidth}p{0.19\linewidth}p{0.48\linewidth}@{}}
\caption{Required array shapes and indexing conventions.}\label{tab:layout}\\
\toprule
Quantity & Julia shape & Meaning \\
\midrule
\endfirsthead
\toprule
Quantity & Julia shape & Meaning \\
\midrule
\endhead
\bottomrule
\endfoot
\texttt{Ypartial} & \texttt{N x M} & $Y_{k,m}$, partial-FFT branch observation for carrier $k$ and branch $m$. \\
\texttt{W} & \texttt{N x M} & $W[k,m]=w_{k,m}$.  Scalar output is $a_k=\mathbf w_k^H\mathbf Y_k$. \\
\texttt{C} & \texttt{N x (2K+1) x M} & $C[k,o,m]=c_{k,q,m}$ with $o=q-k+K+1$. \\
\texttt{z} & \texttt{nbits} & Real bit latents. \\
\texttt{xi} & \texttt{nbits} & $\xi_i=\tanh z_i$. \\
\texttt{tanhprime} & \texttt{nbits} & $r_i=1-\xi_i^2$. \\
\texttt{s} & \texttt{N} & Relaxed carrier symbols.  Data carriers use $\bm z$; pilots are fixed; nulls are zero. \\
\texttt{carrier\_to\_bits} & \texttt{N} entries & For data carrier $q$, returns $(i_I(q),i_Q(q))$.  For pilots/nulls, returns $(0,0)$. \\
\texttt{bit\_to\_carrier} & \texttt{nbits} & Carrier containing bit $i$. \\
\texttt{bit\_to\_axis} & \texttt{nbits} & Either \texttt{:I} or \texttt{:Q}. \\
\end{longtable}

\subsection{Branch-model indexing}

The local ICI neighborhood of carrier $k$ is
\begin{equation}
\mathcal B_K(k)=\{q: q\in\mathcal K,\ |q-k|\le K\}.
\end{equation}
The corresponding offset index is
\begin{equation}
 o(k,q)=q-k+K+1.
\end{equation}
The implementation must skip invalid neighbors:
\begin{equation}
q<1 \quad\text{or}\quad q>N \quad\Longrightarrow\quad \text{skip } q.
\end{equation}
There is no circular wraparound.

\begin{lstlisting}
@inline function neighbor_indices(k::Int, N::Int, K::Int)
    return max(1, k-K):min(N, k+K)
end

@inline function c_offset(k::Int, q::Int, K::Int)::Int
    return q - k + K + 1
end
\end{lstlisting}

\section{Latent bits and relaxed symbols}

The relaxed bipolar bit is
\begin{equation}
\xi_i=\tanh z_i,
\qquad
r_i=\frac{\partial \xi_i}{\partial z_i}=1-\xi_i^2.
\end{equation}
For data carrier $q$, let $(i_I(q),i_Q(q))$ be its two coded-bit indices.  The relaxed QPSK symbol is
\begin{equation}
 s_q(\bm z)=\frac{1}{\sqrt2}\left(\xi_{i_I(q)}+\jj\xi_{i_Q(q)}\right).
\end{equation}
Define the local derivative coefficient
\begin{equation}
\alpha_i=\frac{\partial s_{q(i)}}{\partial \xi_i}
=\begin{cases}
1/\sqrt2, & \text{if bit }i\text{ is the in-phase bit},\\
\jj/\sqrt2, & \text{if bit }i\text{ is the quadrature bit}.
\end{cases}
\end{equation}
Then
\begin{equation}
\frac{\partial s_{q(i)}}{\partial z_i}=\alpha_i(1-\xi_i^2).
\end{equation}

\begin{lstlisting}
@inline function alpha_for_axis(axis::Symbol)::ComplexF64
    invsqrt2 = inv(sqrt(2.0))
    if axis === :I
        return complex(invsqrt2, 0.0)
    elseif axis === :Q
        return complex(0.0, invsqrt2)
    else
        error("axis must be :I or :Q")
    end
end
\end{lstlisting}

\section{Normalized $z$-subproblem}
\label{sec:z-subproblem}

The solver updates $\mathbf C$ and $\mathbf W$ outside the $\bm z$ subproblem.  With $\mathbf C$ and $\mathbf W$ fixed, define
\begin{equation}
\mathbf e_k(\bm z)
=\mathbf Y_k-\sum_{q\in\mathcal B_K(k)}\mathbf c_{kq}s_q(\bm z),
\end{equation}
where $\mathbf Y_k\in\C^M$ and $\mathbf c_{kq}\in\C^M$.

Define the scalar combiner output
\begin{equation}
 a_q=\mathbf w_q^H\mathbf Y_q.
\end{equation}
The implemented normalized loss is
\begin{equation}
\boxed{
\begin{aligned}
\mathcal L_z(\bm z)
&=\frac{\lambda_d}{N_{\rm act}\sigma^2}
\sum_{k\in\mathcal K_{\rm act}}
\norm{\mathbf e_k(\bm z)}_2^2 \\
&\quad +\frac{\rho}{N_D}
\sum_{q\in\mathcal D}
\abs{a_q-s_q(\bm z)}^2 \\
&\quad +\frac{\lambda_p}{N_C}
\sum_{j=1}^{N_C}
\left(1-\prod_{i\in\mathcal N_j}\xi_i\right)^2.
\end{aligned}}
\label{eq:lz}
\end{equation}
Here $N_{\rm act}=|\mathcal K_{\rm act}|$, $N_D=|\mathcal D|$, and $N_C$ is the number of LDPC parity checks.  The normalization makes default weights less dependent on packet size.

\subsection{Default weights}

\begin{table}[h]
\centering
\caption{Default $z$-update constants.}
\begin{tabular}{@{}lll@{}}
\toprule
Parameter & Default & Meaning \\
\midrule
$\lambda_d$ & $1.0$ & branch observation consistency \\
$\rho_{\max}$ & $0.25$ & combiner-symbol tie strength \\
$\lambda_{p,\max}$ & $0.05$ & LDPC parity penalty strength \\
$T_\rho$ & $8$ outer iterations & tie-term ramp time \\
$T_p$ & $12$ outer iterations & parity-term ramp time \\
$z_{\max}$ & $8.0$ & clamp $z_i$ to $[-z_{\max},z_{\max}]$ \\
Gradient clip & $10.0$ & clip $\ell_2$ norm of $\nabla_z\mathcal L_z$ \\
Adam step size & $2\times 10^{-2}$ & default inner update step \\
Adam $(\beta_1,\beta_2,\epsilon)$ & $(0.9,0.999,10^{-8})$ & default Adam constants \\
Finite-difference $\epsilon$ & $10^{-5}$ & gradient-check step \\
Gradient-check tolerance & $10^{-4}$ relative & required finite-difference tolerance \\
\bottomrule
\end{tabular}
\end{table}

The scheduled weights are
\begin{align}
\rho^{(t)}&=\rho_{\max}\left(1-e^{-t/T_\rho}\right)q_{\rm sym}^{(t)},\\
\lambda_p^{(t)}&=\lambda_{p,\max}\left(1-e^{-t/T_p}\right)q_{\rm dec}^{(t)}.
\end{align}
Version 1 uses deterministic gates
\begin{equation}
q_{\rm sym}^{(t)}=q_{\rm dec}^{(t)}=1.
\end{equation}
Confidence-dependent gates may be added later, but they are not part of the first implementation.

\section{Manual gradient for the $z$-subproblem}
\label{sec:z-gradient}

The implementation computes the gradient in two stages:
\begin{enumerate}[label=(\alph*)]
\item accumulate $g_i^{(\xi)}=\partial\mathcal L_z/\partial\xi_i$;
\item apply the tanh chain rule:
\begin{equation}
 g_i^{(z)}=\frac{\partial\mathcal L_z}{\partial z_i}
 =\left(1-\xi_i^2\right)g_i^{(\xi)}.
\end{equation}
\end{enumerate}

\subsection{Observation-consistency gradient}

The observation loss is
\begin{equation}
\mathcal L_d
=\frac{\lambda_d}{N_{\rm act}\sigma^2}
\sum_k \mathbf e_k^H\mathbf e_k.
\end{equation}
For a bit $i$ living on carrier $q(i)$,
\begin{equation}
\frac{\partial \mathcal L_d}{\partial \xi_i}
= -\frac{2\lambda_d}{N_{\rm act}\sigma^2}
\re\left\{
\alpha_i^*\sum_{k:q(i)\in\mathcal B_K(k)}
\mathbf c_{k,q(i)}^H\mathbf e_k
\right\}.
\label{eq:grad_obs}
\end{equation}
This formula follows from
\begin{equation}
\dd\left(\mathbf e^H\mathbf e\right)=2\re\{(\dd\mathbf e)^H\mathbf e\},
\qquad
\frac{\partial\mathbf e_k}{\partial\xi_i}=-\mathbf c_{k,q(i)}\alpha_i.
\end{equation}

\paragraph{Implementation detail.}
Use Julia's \texttt{dot(cvec,e)} for $\mathbf c^H\mathbf e$ because \texttt{dot} conjugates its first argument.

\begin{lstlisting}
function add_observation_gradient!(
    gxi::Vector{Float64},
    e_cache::Matrix{ComplexF64},      # N x M, row k is e_k
    C::Array{ComplexF64,3},           # N x (2K+1) x M
    problem::JUNAProblem,
    cfg::JUNAConfig,
    lambda_d::Float64,
    sigma2::Float64
)::Nothing
    N = problem.ofdm.N
    M = cfg.M
    K = cfg.Kici
    scale = -2.0 * lambda_d / (length(problem.active_indices) * sigma2)

    for k in problem.active_indices
        eview = @view e_cache[k, :]
        for q in neighbor_indices(k, N, K)
            bits = problem.bitmap.carrier_to_bits[q]
            iI, iQ = bits[1], bits[2]
            if iI == 0
                continue  # pilot or null; no z-gradient
            end
            o = c_offset(k, q, K)
            cvec = @view C[k, o, :]
            u = dot(cvec, eview)  # c_{kq}^H e_k

            alphaI = complex(inv(sqrt(2.0)), 0.0)
            alphaQ = complex(0.0, inv(sqrt(2.0)))
            gxi[iI] += scale * real(conj(alphaI) * u)
            gxi[iQ] += scale * real(conj(alphaQ) * u)
        end
    end
    return nothing
end
\end{lstlisting}

\subsection{Combiner-symbol tie gradient}

The tie loss is
\begin{equation}
\mathcal L_\rho
=\frac{\rho}{N_D}\sum_{q\in\mathcal D}\abs{a_q-s_q}^2,
\qquad
 a_q=\mathbf w_q^H\mathbf Y_q.
\end{equation}
Let
\begin{equation}
 r_q=a_q-s_q.
\end{equation}
For bit $i$ on carrier $q(i)$,
\begin{equation}
\frac{\partial \mathcal L_\rho}{\partial\xi_i}
=-\frac{2\rho}{N_D}\re\{\alpha_i^*r_{q(i)}\}.
\label{eq:grad_tie}
\end{equation}

\begin{lstlisting}
function add_tie_gradient!(
    gxi::Vector{Float64},
    shat::Vector{ComplexF64},       # a_q = w_q^H Y_q
    s::Vector{ComplexF64},
    problem::JUNAProblem,
    rho::Float64
)::Nothing
    scale = -2.0 * rho / length(problem.data_indices)
    alphaI = complex(inv(sqrt(2.0)), 0.0)
    alphaQ = complex(0.0, inv(sqrt(2.0)))

    for q in problem.data_indices
        iI, iQ = problem.bitmap.carrier_to_bits[q]
        r = shat[q] - s[q]
        gxi[iI] += scale * real(conj(alphaI) * r)
        gxi[iQ] += scale * real(conj(alphaQ) * r)
    end
    return nothing
end
\end{lstlisting}

\subsection{LDPC parity-penalty gradient}

For check $j$, define
\begin{equation}
 p_j=\prod_{i\in\mathcal N_j}\xi_i,
\qquad
R_j=(1-p_j)^2.
\end{equation}
Then
\begin{equation}
\frac{\partial R_j}{\partial\xi_i}
= -2(1-p_j)\prod_{\ell\in\mathcal N_j\setminus\{i\}}\xi_\ell,
\qquad i\in\mathcal N_j.
\end{equation}
Therefore
\begin{equation}
\frac{\partial\mathcal L_p}{\partial\xi_i}
= -\frac{2\lambda_p}{N_C}
\sum_{j:i\in\mathcal N_j}
(1-p_j)
\prod_{\ell\in\mathcal N_j\setminus\{i\}}\xi_\ell.
\label{eq:grad_parity}
\end{equation}

\paragraph{Important numerical rule.}
Do not compute $\prod_{\ell\ne i}\xi_\ell$ as $p_j/\xi_i$.  If $\xi_i\approx0$, this division is unstable.  Use prefix/suffix products or direct multiplication over the small check neighborhood.

\begin{lstlisting}
function add_parity_gradient!(
    gxi::Vector{Float64},
    xi::Vector{Float64},
    code::LDPCCode,
    lambda_p::Float64
)::Nothing
    nchecks = code.nchecks
    scale = lambda_p / nchecks

    for j in 1:nchecks
        vars = code.check_to_vars[j]
        deg = length(vars)

        # prefix[t] = product xi[vars[1:t-1]]
        prefix = ones(Float64, deg + 1)
        suffix = ones(Float64, deg + 1)

        for t in 1:deg
            prefix[t+1] = prefix[t] * xi[vars[t]]
        end
        for t in deg:-1:1
            suffix[t] = suffix[t+1] * xi[vars[t]]
        end

        pj = prefix[deg+1]
        common = -2.0 * scale * (1.0 - pj)

        for t in 1:deg
            i = vars[t]
            prod_excluding_i = prefix[t] * suffix[t+1]
            gxi[i] += common * prod_excluding_i
        end
    end
    return nothing
end
\end{lstlisting}

\subsection{Chain rule to $z$}

After all $\xi$-gradients have been accumulated,
\begin{equation}
 g_i^{(z)}=(1-\xi_i^2)g_i^{(\xi)}.
\label{eq:chain}
\end{equation}

\begin{lstlisting}
function chain_to_z_gradient!(
    gz::Vector{Float64},
    gxi::Vector{Float64},
    xi::Vector{Float64}
)::Nothing
    @inbounds for i in eachindex(xi)
        gz[i] = (1.0 - xi[i]^2) * gxi[i]
    end
    return nothing
end
\end{lstlisting}

\section{Optional code-aided differential-OFDM gradient}
\label{sec:diff-gradient}

This section is for the optional code-aided differential benchmark.  It is separate from the main SC-JUNA loss.

For a differential pair $(k^-,k)$, define the relaxed differential increment
\begin{equation}
D_k(\bm z)=\frac{1}{\sqrt2}\left(\xi_{i_I(k)}+\jj\xi_{i_Q(k)}\right).
\end{equation}
This is a smooth relaxation of the DQPSK increment alphabet.  The differential residual is
\begin{equation}
 e_k^{\rm diff}=Y_k-D_k(\bm z)Y_{k^-}.
\end{equation}
The normalized differential loss is
\begin{equation}
\mathcal L_{\rm diff}
=\frac{\lambda_{\rm diff}}{N_{\rm diff}}
\sum_{k\in\mathcal D_{\rm diff}}
\omega_k\frac{\abs{e_k^{\rm diff}}^2}{\sigma_{{\rm diff},k}^2}.
\end{equation}
For bit $i$ on differential increment $D_k$,
\begin{equation}
\frac{\partial\mathcal L_{\rm diff}}{\partial\xi_i}
= -\frac{2\lambda_{\rm diff}}{N_{\rm diff}}
\omega_k\frac{1}{\sigma_{{\rm diff},k}^2}
\re\left\{\alpha_i^*Y_{k^-}^*e_k^{\rm diff}\right\}.
\label{eq:grad_diff}
\end{equation}
The same tanh chain rule in \eqref{eq:chain} converts this to a $z$-gradient.

\paragraph{Use case.}
This optional gradient supports a receiver of the form
\begin{equation}
\text{differential metric}+\text{LDPC parity penalty}\rightarrow\text{code-aided differential decoding}.
\end{equation}
It should be implemented as a separate benchmark before being combined with SC-JUNA.

\section{Complete manual-gradient API}

\begin{lstlisting}
struct ManualGradientConfig
    use_observation::Bool
    use_tie::Bool
    use_parity::Bool
    use_differential::Bool
    fd_epsilon::Float64
    fd_rtol::Float64
    grad_clip_norm::Float64
end

struct ZGradientWorkspace
    xi::Vector{Float64}
    tanhprime::Vector{Float64}
    s::Vector{ComplexF64}
    shat::Vector{ComplexF64}
    e_cache::Matrix{ComplexF64}
    gxi::Vector{Float64}
    gz::Vector{Float64}
end
\end{lstlisting}

Required functions:

\begin{lstlisting}
function compute_xi!(ws::ZGradientWorkspace, z::Vector{Float64})::Nothing
function compute_symbols!(ws::ZGradientWorkspace, problem::JUNAProblem)::Nothing
function compute_combiner_outputs!(ws::ZGradientWorkspace, state::JUNAState, obs::RxObservations)::Nothing
function compute_branch_residuals!(ws::ZGradientWorkspace, state::JUNAState, obs::RxObservations, problem::JUNAProblem, cfg::JUNAConfig)::Nothing

function juna_z_loss(
    z::Vector{Float64},
    state::JUNAState,
    obs::RxObservations,
    problem::JUNAProblem,
    cfg::JUNAConfig,
    weights::JUNAIterationWeights,
    ws::ZGradientWorkspace
)::Float64

function juna_z_gradient!(
    gz::Vector{Float64},
    z::Vector{Float64},
    state::JUNAState,
    obs::RxObservations,
    problem::JUNAProblem,
    cfg::JUNAConfig,
    weights::JUNAIterationWeights,
    ws::ZGradientWorkspace
)::Nothing
\end{lstlisting}

\section{Reference implementation pseudocode}

\begin{lstlisting}
function juna_z_gradient!(gz, z, state, obs, problem, cfg, weights, ws)
    # 1. Forward pass
    compute_xi!(ws, z)
    compute_symbols!(ws, problem)
    compute_combiner_outputs!(ws, state, obs)
    compute_branch_residuals!(ws, state, obs, problem, cfg)

    # 2. Accumulate gradient with respect to xi
    fill!(ws.gxi, 0.0)

    if cfg.grad.use_observation
        add_observation_gradient!(
            ws.gxi, ws.e_cache, state.C, problem, cfg,
            weights.lambda_d, obs.noise_power
        )
    end

    if cfg.grad.use_tie
        add_tie_gradient!(ws.gxi, ws.shat, ws.s, problem, weights.rho)
    end

    if cfg.grad.use_parity
        add_parity_gradient!(ws.gxi, ws.xi, problem.ldpc, weights.lambda_p)
    end

    if cfg.grad.use_differential
        add_differential_gradient!(ws.gxi, ws.xi, obs, problem, cfg, weights)
    end

    # 3. Chain rule xi = tanh(z)
    chain_to_z_gradient!(gz, ws.gxi, ws.xi)

    # 4. Optional norm clipping
    clip_gradient_norm!(gz, cfg.grad.grad_clip_norm)

    return nothing
end
\end{lstlisting}

\section{Manual-gradient $z$ update}

Use Adam for the first implementation.  Store Adam moments in the solver state.

\begin{lstlisting}
mutable struct AdamState
    m::Vector{Float64}
    v::Vector{Float64}
    t::Int
end

function adam_step!(z, gz, adam::AdamState; eta=0.02, beta1=0.9, beta2=0.999, eps=1e-8, zmax=8.0)
    adam.t += 1
    @inbounds for i in eachindex(z)
        adam.m[i] = beta1 * adam.m[i] + (1.0 - beta1) * gz[i]
        adam.v[i] = beta2 * adam.v[i] + (1.0 - beta2) * gz[i]^2
        mhat = adam.m[i] / (1.0 - beta1^adam.t)
        vhat = adam.v[i] / (1.0 - beta2^adam.t)
        z[i] -= eta * mhat / (sqrt(vhat) + eps)
        z[i] = clamp(z[i], -zmax, zmax)
    end
    return nothing
end
\end{lstlisting}

\section{Closed-form updates for $\mathbf C$ and $\mathbf W$}

No gradient or AD is needed for $\mathbf C$ or $\mathbf W$ in version 1.

\subsection{$\mathbf C$ update}

For each carrier $k$, let
\begin{equation}
\mathbf s_k=[s_q]_{q\in\mathcal B_K(k)}\in\C^{B_k},
\qquad
\mathbf C_k=[\mathbf c_{kq_1},\ldots,\mathbf c_{kq_{B_k}}]\in\C^{M\times B_k}.
\end{equation}
With diagonal noise weight, the ridge subproblem is
\begin{equation}
\min_{\mathbf C_k}
\frac{\lambda_d}{N_{\rm act}\sigma^2}
\norm{\mathbf Y_k-\mathbf C_k\mathbf s_k}_2^2
+\gamma_C\norm{\mathbf C_k}_F^2.
\end{equation}
Define
\begin{equation}
\mu_C=\frac{\lambda_d}{N_{\rm act}\sigma^2}.
\end{equation}
The solution is
\begin{equation}
\boxed{
\mathbf C_k
=\mu_C\mathbf Y_k\mathbf s_k^H
\left(\mu_C\mathbf s_k\mathbf s_k^H+\gamma_C I\right)^{-1}.}
\end{equation}

\subsection{$\mathbf W$ update without smoothness}

For each carrier $k$, set
\begin{equation}
d_k=\begin{cases}
P_k, & k\in\mathcal P_{\rm fit},\\
s_k, & k\in\mathcal D,
\end{cases}
\qquad
a_k=\begin{cases}
\eta/N_P, & k\in\mathcal P_{\rm fit},\\
\rho/N_D, & k\in\mathcal D.
\end{cases}
\end{equation}
The ridge subproblem is
\begin{equation}
\min_{\mathbf w_k}
a_k\abs{\mathbf w_k^H\mathbf Y_k-d_k}^2+\gamma_W\norm{\mathbf w_k}_2^2.
\end{equation}
The solution is
\begin{equation}
\boxed{
\mathbf w_k=
\left(a_k\mathbf Y_k\mathbf Y_k^H+\gamma_W I\right)^{-1}
a_k\mathbf Y_k d_k^*.}
\end{equation}

\subsection{$\mathbf W$ update with smoothness}

If smoothness is enabled, solve
\begin{equation}
\sum_k a_k\abs{\mathbf w_k^H\mathbf Y_k-d_k}^2
+\gamma_W\sum_k\norm{\mathbf w_k}_2^2
+\gamma_S\sum_k\norm{\mathbf w_{k+1}-\mathbf w_k}_2^2.
\end{equation}
This gives a block-tridiagonal linear system.  Interior block equations are
\begin{equation}
\left(a_k\mathbf Y_k\mathbf Y_k^H+(\gamma_W+2\gamma_S)I\right)\mathbf w_k
-\gamma_S\mathbf w_{k-1}
-\gamma_S\mathbf w_{k+1}
=a_k\mathbf Y_k d_k^*.
\end{equation}
Endpoint carriers use one smoothness neighbor instead of two.  Version 1 may set $\gamma_S=0$ and use the independent formula above.  Version 2 should implement the block-tridiagonal solve.

\section{Gradient checks}
\label{sec:checks}

Every gradient component must pass a finite-difference test before integration.

\subsection{Central finite difference}

For a random test direction $\bm v$ with $\norm{\bm v}_2=1$,
\begin{equation}
D_{\rm fd}=\frac{\mathcal L_z(\bm z+\epsilon\bm v)-\mathcal L_z(\bm z-\epsilon\bm v)}{2\epsilon}.
\end{equation}
The analytic directional derivative is
\begin{equation}
D_{\rm an}=\bm v^T\nabla_z\mathcal L_z(\bm z).
\end{equation}
Pass if
\begin{equation}
\frac{\abs{D_{\rm fd}-D_{\rm an}}}{\max(1,\abs{D_{\rm fd}},\abs{D_{\rm an}})}<10^{-4}.
\end{equation}

\begin{lstlisting}
function gradient_check_directional(
    lossfun::Function,
    gradfun!::Function,
    z::Vector{Float64};
    eps::Float64 = 1e-5,
    rtol::Float64 = 1e-4,
    rng::AbstractRNG = Random.default_rng()
)::Bool
    v = randn(rng, length(z))
    v ./= norm(v)

    gz = similar(z)
    gradfun!(gz, z)

    fd = (lossfun(z .+ eps .* v) - lossfun(z .- eps .* v)) / (2eps)
    an = dot(v, gz)
    relerr = abs(fd - an) / max(1.0, abs(fd), abs(an))
    return relerr < rtol
end
\end{lstlisting}

\subsection{Required unit tests}

\begin{enumerate}[label=\textbf{G\arabic*.},leftmargin=*]
\item \textbf{Tie-only gradient.}  Set $\lambda_d=0$, $\lambda_p=0$, $\rho>0$.  Compare analytic gradient against finite difference.
\item \textbf{Observation-only gradient.}  Set $\rho=0$, $\lambda_p=0$, $\lambda_d>0$.  Use $N=5$, $M=3$, $K=1$, random fixed $\mathbf C$, and finite-difference check.
\item \textbf{Parity-only gradient.}  Use one parity check $\mathcal N_1=\{1,2,3\}$.  Compare analytic gradient against finite difference.
\item \textbf{Full $z$ gradient.}  Enable observation, tie, and parity terms.  Test ten random directions.
\item \textbf{Differential gradient.}  Enable only $\mathcal L_{\rm diff}$ for a small synthetic differential chain.  Test finite difference.
\item \textbf{No invalid edge access.}  Use $K=3$ and $N=4$.  Confirm all neighbor loops skip invalid carriers and do not wrap.
\end{enumerate}

\section{Solver loop using manual gradients}

\begin{lstlisting}
function run_juna_solver_manualgrad(
    state::JUNAState,
    obs::RxObservations,
    problem::JUNAProblem,
    cfg::JUNAConfig,
    bp_cfg::BPConfig
)::JUNAResult
    ws = allocate_z_gradient_workspace(state, obs, problem, cfg)
    adam = AdamState(zeros(length(state.z)), zeros(length(state.z)), 0)
    gz = zeros(length(state.z))

    for outer in 1:cfg.max_outer
        weights = iteration_weights(cfg, outer)

        # 1. Update relaxed symbols from current z
        compute_xi!(ws, state.z)
        compute_symbols!(ws, problem)
        copyto!(state.s, ws.s)

        # 2. Closed-form front-end updates
        update_C_closed_form!(state, obs, problem, cfg, weights)
        update_W_closed_form!(state, obs, problem, cfg, weights)

        # 3. Manual-gradient z update
        for inner in 1:cfg.z_steps_per_outer
            juna_z_gradient!(gz, state.z, state, obs, problem, cfg, weights, ws)
            adam_step!(
                state.z, gz, adam;
                eta = cfg.z_stepsize,
                beta1 = 0.9,
                beta2 = 0.999,
                eps = 1e-8,
                zmax = cfg.z_clip
            )
        end

        # 4. Form LLRs, run BP, stop on parity success
        llr = form_llrs_from_state(state, obs, problem, cfg)
        bp = bp_decode(llr, problem.ldpc, bp_cfg)
        push!(state.loss_history, juna_z_loss(state.z, state, obs, problem, cfg, weights, ws))
        push!(state.syndrome_history, bp.syndrome_weight)

        if bp.success
            return JUNAResult(bp.bits, llr, bp, state, cfg, nothing, obs.diagnostics)
        end
    end

    llr = form_llrs_from_state(state, obs, problem, cfg)
    bp = bp_decode(llr, problem.ldpc, bp_cfg)
    return JUNAResult(bp.bits, llr, bp, state, cfg, nothing, obs.diagnostics)
end
\end{lstlisting}

\section{LLR reliability and erasure}

Manual gradients do not replace reliability control.  After the front-end produces scalar soft symbols
\begin{equation}
\widehat s_k=\mathbf w_k^H\mathbf Y_k,
\end{equation}
compute an effective variance from pilot residuals:
\begin{equation}
\widehat\sigma_{{\rm eff},k}^2
=\operatorname{localmedian}_{p\approx k}
\abs{\mathbf w_p^H\mathbf Y_p-P_p}^2+\epsilon.
\end{equation}
Then QPSK LLRs are
\begin{align}
L_{I,k}&=\frac{2\sqrt2\re\{\widehat s_k\}}{\widehat\sigma_{{\rm eff},k}^2},\\
L_{Q,k}&=\frac{2\sqrt2\im\{\widehat s_k\}}{\widehat\sigma_{{\rm eff},k}^2}.
\end{align}
Clip and erase:
\begin{align}
L_i&\leftarrow \clip(L_i,-L_{\max},L_{\max}),\\
\gamma_k&=\frac{1}{\widehat\sigma_{{\rm eff},k}^2+\epsilon},\\
\gamma_k<\tau_{\rm erase}&\Longrightarrow L_{I,k}=L_{Q,k}=0.
\end{align}
Defaults:
\begin{equation}
L_{\max}=12,
\qquad
\tau_{\rm erase}=0.05\cdot\operatorname{median}_k(\gamma_k).
\end{equation}

\section{When to use AD}

AD is not needed for the production solver.  Acceptable uses of AD are:
\begin{enumerate}[label=(\alph*)]
\item checking the manual gradient in small real-valued toy problems;
\item prototyping new penalty terms before deriving their gradients;
\item validating a future non-diagonal $\Sigma_k^{-1}$ implementation.
\end{enumerate}

If AD is used for checking, the loss must be a real scalar function of a real vector:
\begin{equation}
\mathcal L:\R^{n_b}\rightarrow\R.
\end{equation}
Do not expose complex parameters as AD variables in the first implementation.  Keep $\mathbf C$, $\mathbf W$, and observations fixed while checking $\nabla_z\mathcal L_z$.

\section{Implementation order}

\begin{enumerate}[label=\textbf{Step \arabic*.},leftmargin=*]
\item Implement array-layout helpers: neighbor loops, offset mapping, carrier-bit map, and axis lookup.
\item Implement forward pass: $\bm z\mapsto\bm\xi\mapsto\mathbf s$, combiner outputs, branch residuals.
\item Implement scalar loss terms: observation loss, tie loss, parity loss, optional differential loss.
\item Implement manual gradient components: observation, tie, parity, optional differential.
\item Implement finite-difference tests for each component.
\item Implement $\mathbf C$ closed-form update.
\item Implement $\mathbf W$ closed-form update without smoothness.
\item Integrate Adam $z$ update.
\item Add BP decoder call and stopping rule.
\item Add SC-JUNA candidate selection only after the fixed JUNA manual-gradient solver passes all tests.
\item Add smooth $\mathbf W$ block-tridiagonal solve as a version-2 improvement.
\item Add optional AD gradient checker only after the finite-difference checker passes.
\end{enumerate}

\section{Minimum acceptance checklist}

The implementation is not acceptable until the following pass:
\begin{enumerate}[leftmargin=*]
\item Partial-FFT partition-of-unity test passes.
\item LDPC noiseless BP decode passes.
\item Tie-only gradient finite-difference test passes.
\item Observation-only gradient finite-difference test passes.
\item Parity-only gradient finite-difference test passes.
\item Full $z$ gradient finite-difference test passes on ten random directions.
\item $\mathbf C$ update reduces branch observation residual in a deterministic toy case.
\item $\mathbf W$ update reduces pilot residual in a deterministic toy case.
\item Fixed JUNA runs without AD and produces finite loss values for all iterations.
\item LLR clipping and erasure do not produce NaN or Inf.
\end{enumerate}

\section{Summary}

The hand-derived gradient is straightforward because all optimized bit variables are real.  The complex channel and branch quantities remain fixed during the $\bm z$ update, so the gradient only requires real derivatives through the relaxed QPSK map and the LDPC parity products.  The recommended implementation is:
\begin{equation}
\boxed{
\text{closed-form }\mathbf C
+\text{ closed-form }\mathbf W
+\text{ manual }\nabla_z
+\text{ finite-difference gradient checks}.}
\end{equation}
AD can be useful for experimentation, but it should not be the production dependency for the SC-JUNA solver.

\end{document}
